\section{Definitions of Evaluation Metrics}
\label{appendix:metrics}

Let $P$ and $G$ denote the sets of predicted and reference lesion voxels of
a test case, and $|\cdot|$ the number of voxels in a set. Let
$\mathcal{G}$ and $\mathcal{P}$ denote the sets of reference and predicted
lesions (6-connected components), and
$\mathcal{G}_{\mathrm{det}} \subseteq \mathcal{G}$ and
$\mathcal{P}_{\mathrm{det}} \subseteq \mathcal{P}$ the reference lesions
that overlap at least one predicted lesion and the predicted lesions that
overlap at least one reference lesion. In all formulas,
$\varepsilon = 10^{-8}$ prevents division by zero. Every metric is computed
per test case and then averaged over cases.

\paragraph{Voxel-level metrics.}
\begin{equation}
\mathrm{DSC} = \frac{2|P \cap G|}{|P| + |G|}.
\end{equation}
With $S_P$ and $S_G$ the surface voxels of the two masks and
$d(p, S_G) = \min_{g \in S_G} \lVert p - g \rVert_2$ in mm,
\begin{equation}
\mathrm{HD}_{95}(P,G) =
\max \left\{
\underset{p \in S_P}{\mathrm{percentile}_{95}}\, d(p, S_G),\;
\underset{g \in S_G}{\mathrm{percentile}_{95}}\, d(g, S_P)
\right\}.
\end{equation}

\paragraph{Lesion-level metrics.}
\begin{equation}
\mathrm{Precision}_{\mathrm{les}} =
\frac{|\mathcal{P}_{\mathrm{det}}|}{|\mathcal{P}| + \varepsilon},
\qquad
\mathrm{Recall}_{\mathrm{les}} =
\frac{|\mathcal{G}_{\mathrm{det}}|}{|\mathcal{G}| + \varepsilon},
\qquad
\mathrm{F1}_{\mathrm{les}} =
\frac{2 \, \mathrm{Precision}_{\mathrm{les}} \, \mathrm{Recall}_{\mathrm{les}}}
{\mathrm{Precision}_{\mathrm{les}} + \mathrm{Recall}_{\mathrm{les}} + \varepsilon}.
\end{equation}
With $\tau = 150$ voxels and
$\mathcal{G}_{\mathrm{small}} = \{ g \in \mathcal{G} : |g| \leq \tau \}$,
\begin{equation}
\mathrm{Recall}_{\mathrm{small}} =
\frac{|\mathcal{G}_{\mathrm{small}} \cap \mathcal{G}_{\mathrm{det}}|}
{|\mathcal{G}_{\mathrm{small}}| + \varepsilon},
\qquad
\text{FN Count} = |\mathcal{G}| - |\mathcal{G}_{\mathrm{det}}|.
\end{equation}
The recall of lesions of at least 3~mm is $\mathrm{Recall}_{\mathrm{les}}$
restricted to reference lesions whose longest diameter is at least 3~mm.
The longest diameter of a lesion is the maximum, over all axial, coronal,
and sagittal slices that intersect it, of the largest distance between two
of its voxel centers in that slice, plus one voxel width (the mean in-plane
spacing), so that it is measured from edge to edge.

\paragraph{Error-volume metrics.}
Let $\mathrm{FP} = P \setminus G$ denote the false-positive voxels and
$s_x, s_y, s_z$ the voxel spacing. The false-positive volume is
\begin{equation}
V_{\mathrm{FP}} = |\mathrm{FP}| \cdot s_x s_y s_z.
\end{equation}
The FN Volume Fraction is
\begin{equation}
V_{\mathrm{FN}}^{\mathrm{frac}}
= \frac{\sum_{g \in \mathcal{G} \setminus \mathcal{G}_{\mathrm{det}}} |g|}
{|G| + \varepsilon}.
\end{equation}

\section{Implementation and Statistical Details}
\label{appendix:impl-stats}

\paragraph{Implementation.}
\label{appendix:implementation}
All models use nnU-Net (version 2.1.1)~\cite{isensee_nnu-net_2021} with
its default optimizer: stochastic gradient descent with Nesterov momentum
of 0.99, weight decay of $3\times10^{-5}$, and an initial learning rate of
0.01 with polynomial decay. Training runs for 150 epochs with a batch size
of 2 on an NVIDIA Quadro RTX 8000 GPU. Tversky and FocalTversky replace the
DiceCE loss, whereas CC-DiceCE and CATMIL add terms to it. CC-DiceCE adds
its per-lesion term, built on the component-wise evaluation protocol of
CC-Metrics~\cite{jaus_every_2025}, to the global DiceCE loss with equal
weight. FocalTversky is $\mathcal{L}_{\mathrm{Tversky}}^{1/\gamma_{\mathrm{F}}}$
with the same $\alpha$ and $\beta$ as Tversky. The auxiliary terms of CATMIL
introduce no parameters and are used only during training, so the inference
cost is unchanged; their training-time cost is dominated by connected-component
labelling of the ground-truth patches at each iteration.

\paragraph{Statistical analysis.}
\label{appendix:stats}
For each metric, the difference between CATMIL and a baseline is computed
in the same way as the table values: the mean over test cases of the
five-model mean. Its 95\% CI is obtained with a two-level bootstrap with
10{,}000 resamples. Each resample draws patients with replacement, together
with all their scans, and draws fold models with replacement. The fold
models are drawn jointly for both losses, because fold $k$ of every loss is
trained on the same data split. For lesion-level recall, the same bootstrap
is applied to the pooled reference lesions of the resampled patients. A
paired exact Wilcoxon signed-rank test is applied to the per-patient
differences, with the Holm correction over the four baselines. With six
patients, the smallest possible two-sided $p$-value of the exact test is
$1/32 \approx 0.031$ before correction and $0.125$ after it.

\section{Full Paired Differences on the Test Sets}
\label{appendix:paired-full}

\Cref{tab:paired_differences_full} extends \cref{tab:paired_differences}
to all metrics and to the MSLesSeg results after PP5.

\begin{table}[t]
\centering
\caption{All paired differences between CATMIL and each baseline on the test sets, formatted as in
\cref{tab:paired_differences}. With six patients, the smallest possible
Holm-corrected Wilcoxon $p$-value is 0.125 (\cref{sec:stats}).}
\label{tab:paired_differences_full}
\resizebox{\textwidth}{!}{%
\begin{tabular}{lllll}
\toprule
\textbf{Metric} & \textbf{vs.\ DiceCE} & \textbf{vs.\ Tversky} & \textbf{vs.\ FocalTversky} & \textbf{vs.\ CC-DiceCE} \\
\midrule
\multicolumn{5}{l}{\emph{MSLesSeg}} \\
Small Lesion Recall $\uparrow$ & \textbf{+0.077 [+0.030, +0.157] (6/6)} & \textbf{+0.039 [+0.004, +0.093] (6/6)} & \textbf{+0.042 [+0.005, +0.101] (6/6)} & \textbf{+0.042 [+0.020, +0.071] (6/6)} \\
FN Count $\downarrow$ & \textbf{$-$2.38 [$-$5.29, $-$0.84] (6/6)} & \textbf{$-$1.20 [$-$2.98, $-$0.27] (6/6)} & \textbf{$-$1.13 [$-$2.71, $-$0.27] (5/6)} & \textbf{$-$1.08 [$-$1.95, $-$0.52] (6/6)} \\
DSC $\uparrow$ & +0.004 [$-$0.010, +0.031] (1/6) & \textbf{+0.013 [+0.001, +0.034] (5/6)} & +0.003 [$-$0.015, +0.023] (5/6) & +0.006 [$-$0.003, +0.017] (6/6) \\
HD95 (mm) $\downarrow$ & $-$1.06 [$-$4.67, +0.73] (4/6) & $-$2.26 [$-$7.35, +0.11] (4/6) & $-$0.22 [$-$2.92, +2.72] (4/6) & \textbf{$-$1.73 [$-$5.43, $-$0.01] (5/6)} \\
Lesion F1 $\uparrow$ & \textbf{$-$0.086 [$-$0.138, $-$0.043] (0/6)} & \textbf{$-$0.083 [$-$0.143, $-$0.038] (0/6)} & \textbf{$-$0.088 [$-$0.151, $-$0.039] (0/6)} & \textbf{$-$0.049 [$-$0.091, $-$0.018] (0/6)} \\
FP Volume (mm$^3$) $\downarrow$ & $-$282 [$-$1582, +271] (1/6) & \textbf{$-$1084 [$-$2675, $-$302] (6/6)} & \textbf{$-$746 [$-$2189, $-$182] (6/6)} & +23 [$-$276, +236] (1/6) \\
\midrule
\multicolumn{5}{l}{\emph{MSLesSeg, PP5}} \\
Small Lesion Recall $\uparrow$ & \textbf{+0.070 [+0.023, +0.151] (6/6)} & \textbf{+0.033 [+0.001, +0.079] (6/6)} & \textbf{+0.036 [+0.000, +0.093] (6/6)} & \textbf{+0.040 [+0.017, +0.073] (6/6)} \\
FN Count $\downarrow$ & \textbf{$-$2.10 [$-$4.80, $-$0.66] (6/6)} & \textbf{$-$0.93 [$-$2.38, $-$0.13] (6/6)} & \textbf{$-$0.95 [$-$2.40, $-$0.14] (6/6)} & \textbf{$-$1.05 [$-$1.84, $-$0.46] (6/6)} \\
DSC $\uparrow$ & +0.004 [$-$0.009, +0.032] (2/6) & \textbf{+0.013 [+0.001, +0.035] (5/6)} & +0.004 [$-$0.014, +0.024] (5/6) & +0.006 [$-$0.003, +0.017] (6/6) \\
HD95 (mm) $\downarrow$ & $-$0.99 [$-$4.65, +0.76] (4/6) & $-$2.26 [$-$7.25, +0.10] (5/6) & $-$0.13 [$-$2.80, +2.76] (4/6) & \textbf{$-$1.72 [$-$5.34, $-$0.11] (5/6)} \\
Lesion F1 $\uparrow$ & $-$0.017 [$-$0.040, +0.001] (2/6) & $-$0.012 [$-$0.038, +0.006] (2/6) & $-$0.017 [$-$0.043, +0.001] (2/6) & +0.003 [$-$0.010, +0.017] (4/6) \\
FP Volume (mm$^3$) $\downarrow$ & $-$301 [$-$1610, +244] (1/6) & \textbf{$-$1103 [$-$2706, $-$310] (6/6)} & \textbf{$-$765 [$-$2211, $-$197] (6/6)} & +9 [$-$295, +216] (1/6) \\
\midrule
\multicolumn{5}{l}{\emph{3D-MR-MS}} \\
Small Lesion Recall $\uparrow$ & \textbf{+0.081 [+0.034, +0.127] (5/6)} & \textbf{+0.054 [+0.019, +0.086] (5/6)} & \textbf{+0.055 [+0.021, +0.087] (5/6)} & +0.039 [$-$0.001, +0.080] (4/6) \\
FN Count $\downarrow$ & \textbf{$-$5.13 [$-$12.40, $-$0.50] (5/6)} & \textbf{$-$4.43 [$-$11.23, $-$0.27] (5/6)} & \textbf{$-$4.43 [$-$11.47, $-$0.23] (5/6)} & \textbf{$-$2.27 [$-$4.80, $-$0.47] (5/6)} \\
DSC $\uparrow$ & $-$0.012 [$-$0.032, +0.010] (0/6) & $-$0.008 [$-$0.030, +0.013] (3/6) & $-$0.016 [$-$0.035, +0.001] (1/6) & $-$0.002 [$-$0.015, +0.013] (2/6) \\
HD95 (mm) $\downarrow$ & $-$0.38 [$-$4.14, +2.07] (3/6) & $-$1.83 [$-$8.51, +1.50] (2/6) & $-$0.72 [$-$6.07, +2.92] (3/6) & $-$2.36 [$-$8.87, +0.60] (4/6) \\
Lesion F1 $\uparrow$ & \textbf{$-$0.197 [$-$0.258, $-$0.142] (0/6)} & \textbf{$-$0.168 [$-$0.242, $-$0.106] (0/6)} & \textbf{$-$0.166 [$-$0.232, $-$0.108] (0/6)} & \textbf{$-$0.131 [$-$0.195, $-$0.066] (0/6)} \\
FP Volume (mm$^3$) $\downarrow$ & $-$2 [$-$110, +110] (3/6) & \textbf{$-$761 [$-$1587, $-$118] (6/6)} & \textbf{$-$726 [$-$1550, $-$92] (6/6)} & $-$50 [$-$145, +16] (3/6) \\
\bottomrule
\end{tabular}%
}
\end{table}

\section{Ablation I: Contribution of CAT and MIL Terms}
\label{appendix:catmil-components}

The ablation uses the nnU-Net configuration, data splits, and training
schedule of the main experiments and compares three variants: (i) CAT,
which adds only the component-adaptive term to the base loss, (ii) MIL,
which adds only the lesion-level term, and (iii) CATMIL, which combines
both.

\Cref{tab:ablation_cat_mil} reports the mean performance of the five fold
models on the MSLesSeg test set, and the differences below are given with
95\% patient-level CIs as in \cref{tab:paired_differences}.

\textbf{The MIL term accounts for the detection gain.} MIL alone raises
small-lesion recall from 0.7956 (DiceCE) to 0.8712 (+0.076 [+0.033,
+0.148]) and reduces missed lesions per case by 2.42 [0.87, 5.26], in all
six patients. These gains are as large as those of CATMIL, and CATMIL and
MIL do not differ on any metric: small-lesion recall differs by +0.002
[$-0.012$, +0.019], lesion-wise F1 by +0.008 [$-0.020$, +0.031], DSC by
+0.003 [$-0.016$, +0.021], and HD95 by $-1.26$~mm [$-6.82$, +1.55]. The
lower lesion-wise F1 of CATMIL is also already present with MIL alone
(0.7492).

\textbf{The CAT term alone changes lesion detection little.} CAT alone
reaches a small-lesion recall of 0.8230, a difference from DiceCE of
+0.027 [$-0.010$, +0.081] whose CI includes zero, and it keeps the
lesion-wise F1 of DiceCE (0.8386 against 0.8433). Its effect on detection
is thus close to that of the Tversky-based losses and of CC-DiceCE
(0.8307), which uses the same per-lesion weighting: the term rebalances
voxel weights between lesions but, unlike the MIL term, does not require
each lesion to be detected.

The mean DSC and HD95 of CATMIL are the best of the three variants, and
CATMIL has the lowest false-positive blob fraction, but these differences
lie within the variation between patients and fold models. On this test
set, the ablation therefore supports the MIL term as the source of the
detection gain, and it does not show a measurable benefit of adding the
CAT term to it. The CAT term is kept in CATMIL because of its design
rationale, size-balanced voxel supervision within each lesion
(\cref{appendix:catmil-properties}), and because the loss-weight analysis
below shows that its weight shifts the balance between voxel-level accuracy
and lesion-wise precision; whether it contributes measurably on other data
remains to be tested.

\begin{table}[H]
\centering
\caption{Ablation study on the individual and combined contribution of the
CAT and MIL terms on the MSLesSeg test set. FP Vol: false-positive volume
(mm$^3$); FP Blob: false-positive blob fraction.}
\label{tab:ablation_cat_mil}
\begin{tabular}{lcccccc}
\toprule
Model & DSC & HD95 $\downarrow$ & Lesion F1 & Small Lesion Recall & FP Vol $\downarrow$ & FP Blob $\downarrow$ \\
\midrule
CAT    & 0.7812 & 8.6844 & \textbf{0.8386} & 0.8230 & \textbf{1528.72} & 0.2277 \\
MIL    & 0.7805 & 9.2375 & 0.7492 & 0.8712 & 1539.77 & 0.2195 \\
CATMIL & \textbf{0.7834} & \textbf{7.9817} & 0.7571 & \textbf{0.8730} & 1536.88 & \textbf{0.2094} \\
\bottomrule
\end{tabular}
\end{table}

\section{Ablation II: Sensitivity to Loss Weights}
\label{appendix:catmil-loss-weights}

This section analyzes the impact of the weighting coefficients
$\lambda_{\text{CAT}}$ and $\lambda_{\text{MIL}}$ (\cref{tab:ablation_lambda}). The configurations are
named CATMIL$xy$, where $\lambda_{\text{CAT}}=0.x$ and
$\lambda_{\text{MIL}}=0.y$; CATMIL0102 is the configuration used in the
main experiments, whose weights were fixed beforehand; the analysis is a
post hoc sensitivity check, not a selection. Each configuration was trained
once on the complete development set, without cross-validation, and
evaluated on the same test set. The values of CATMIL0102 therefore differ from those of
CATMIL in \cref{tab:loss_comparison_small_lesion,tab:ablation_cat_mil},
which are means over the five cross-validation models.

\begin{table}[H]
\centering
\caption{Sensitivity of the CATMIL objective to $\lambda_{\text{CAT}}$ and
$\lambda_{\text{MIL}}$ on the MSLesSeg test set.}
\label{tab:ablation_lambda}
\begin{tabular}{lccccc}
\toprule
Model & DSC & HD95 $\downarrow$ & Lesion F1 & Small Recall & FP Vol $\downarrow$ \\
\midrule
CATMIL0101 & 0.7675 & 9.3881 & \underline{0.7720} & 0.8514 & 2292.17 \\
CATMIL0102 & \textbf{0.7953} & \textbf{5.9771} & 0.7452 & \underline{0.8711} & \underline{1345.83} \\
CATMIL0201 & \underline{0.7921} & \underline{6.7964} & \textbf{0.7754} & \textbf{0.8784} & \textbf{1233.92} \\
CATMIL0202 & 0.7663 & 9.4506 & 0.7632 & 0.8699 & 1729.42 \\
\bottomrule
\end{tabular}
\end{table}

The configuration used in the main experiments (CATMIL0102), in which the
MIL term has the larger weight, gives the best global segmentation
performance, with the highest DSC (0.7953) and lowest HD95 (5.9771~mm), but
also the lowest lesion-wise F1 (0.7452). Exchanging the two weights
(CATMIL0201) improves lesion-level behavior, with the highest lesion F1
(0.7754), the highest small lesion recall (0.8784), and the lowest FP
volume (1233.92~mm$^3$), at the cost of slightly lower DSC and boundary
accuracy. The two configurations with equal weights (CATMIL0101 and
CATMIL0202) are clearly worse in DSC, HD95, and FP volume. Small lesion
recall stays between 0.85 and 0.88 in all four configurations, so the
detection gain does not depend on a precise choice of the weights in this
range; the ratio of $\lambda_{\text{CAT}}$ to $\lambda_{\text{MIL}}$ mainly
shifts the balance between voxel-level accuracy and lesion-wise precision.

\section{Lesion-Profile Characteristics of Test Cases}
\label{appendix:test-case-lesion-profile}

Each MSLesSeg test case is summarized using the ground-truth lesion mask
(\cref{tab:test_case_lesion_profile}). Lesions are the 6-connected
components of the binary mask, as in the loss and the metrics, counted
regardless of size. Lesion volume is the voxel count times the voxel volume
(1~mm$^3$ in MSLesSeg); the total and median volumes are taken over the
lesions of one case.

Lesion contrast is computed locally on the FLAIR image after z-score
normalization within the brain. For each lesion, a surrounding background
ring is created within 3~mm around the lesion, excluding the lesion itself.
The contrast value is the absolute difference between the median intensity
inside the lesion and the median intensity in this ring, divided by the
median absolute deviation of the ring intensity.

Boundary sharpness is computed from the image gradient around the lesion
boundary. The normalized FLAIR image is smoothed with a Gaussian filter of
1~mm, and the sharpness is the median gradient magnitude in a band of one
voxel on each side of the lesion boundary.

Contrast and sharpness are not estimated reliably for the smallest lesions
and are therefore computed only for lesions of at least 10 voxels. The
contrast and sharpness in \cref{tab:test_case_lesion_profile} are means
over these lesions, whereas the lesion count and the volume statistics
include all lesions.

\begin{table}[H]
\centering
\caption{Lesion-profile characteristics of the MSLesSeg test cases. N:
number of lesions; $\leq$50 and $\leq$100: percentage of lesions of at most
50 and 100~mm$^3$.}
\label{tab:test_case_lesion_profile}
\begin{tabular}{lrrrrrrr}
\toprule
\textbf{Case} & \textbf{N} & \textbf{Total vol.} & \textbf{Median vol.}
& \textbf{$\leq$50} & \textbf{$\leq$100} & \textbf{Contrast} & \textbf{Sharpness} \\
& & \textbf{(mm$^3$)} & \textbf{(mm$^3$)} & \textbf{(\%)} & \textbf{(\%)} & & \\
\midrule
P13\_T1 & 39 & 9677 & 83.0 & 10.3 & 53.8 & 4.11 & 8.41 \\
P13\_T2 & 61 & 10268 & 68.0 & 45.9 & 65.6 & 3.85 & 8.34 \\
P26\_T1 & 43 & 6639 & 106.0 & 20.9 & 46.5 & 3.77 & 9.16 \\
P2\_T1  & 16 & 2261 & 92.5 & 31.2 & 50.0 & 4.56 & 9.31 \\
P2\_T2  & 20 & 3286 & 87.5 & 20.0 & 55.0 & 3.76 & 10.30 \\
P2\_T3  & 14 & 2623 & 134.0 & 14.3 & 28.6 & 4.61 & 11.68 \\
P2\_T4  & 18 & 2940 & 87.0 & 16.7 & 55.6 & 4.45 & 10.86 \\
P47\_T1 & 59 & 18736 & 67.0 & 42.4 & 71.2 & 4.26 & 9.92 \\
P49\_T1 & 30 & 34709 & 83.5 & 30.0 & 56.7 & 3.65 & 8.55 \\
P49\_T2 & 19 & 7763 & 173.0 & 0.0 & 26.3 & 3.34 & 7.38 \\
P52\_T1 & 28 & 3306 & 70.5 & 42.9 & 67.9 & 5.24 & 11.02 \\
P52\_T2 & 14 & 2843 & 147.0 & 7.1 & 42.9 & 4.04 & 6.91 \\
\bottomrule
\end{tabular}
\end{table}

\section{Selection of the Lesion Examples}
\label{appendix:gallery-rule}

The lesions in \cref{fig:qual_lesion_gallery} were selected by a fixed rule
applied to the detection results of all five fold models, not by visual
inspection. A \emph{recovered} lesion has at most 50 voxels, is detected by
at least four of the five CATMIL models, and by at most one of the five
models of each voxel-level baseline; a lesion \emph{missed by all} is
detected by none of the twenty models; the \emph{CATMIL false positive} is
a component predicted by all five CATMIL models and by none of the baseline
models. CC-DiceCE was not part of this case analysis.

\section{Additional Lesion-Level Analyses}
\label{appendix:lesion-level-analyses}

These analyses complement the size analysis in the main text and use the
lesion definition of the evaluation metrics: 6-connected components of the reference mask, with
a lesion counted as detected if any predicted voxel overlaps it. Values are
means over the five fold models.

\paragraph{Operating points under component filtering.}
PP5 uses a single filter size. \Cref{fig:operating_curve} shows how lesion
recall and the number of false-positive components per case change when
the filter size is varied. In the left panel, predicted components smaller
than $s$ voxels are removed, with $s$ from 0 to 20; in the right panel, the
probability threshold $t$ is varied from 0.1 to 0.9 without filtering.

The CATMIL curve lies above the baseline curves across the whole range of
filter sizes. Without filtering, CATMIL produces more false-positive
components than the baselines (20.3 per case, against 5.4--7.0 for the
voxel-level losses and 11.2 for CC-DiceCE), which is the cause of its lower
Lesion F1. These components are small and are
removed at little cost in recall. With $s = 20$, CATMIL detects 84.5\% of
lesions at 3.3 false-positive components per case, more lesions with fewer
false alarms than unfiltered DiceCE (83.5\% at 5.4). For any false-alarm
budget of at least 3.3 components per case, a filtered CATMIL model
therefore detects more lesions than any of the baselines. Raising the
probability threshold is less effective than a size filter, because most
of the false-positive components are confident predictions rather than
borderline ones. The margin over the Tversky-based losses at a fixed
false-alarm level is small, however, and is not supported by the
patient-level analysis at matched precision (\cref{tab:matched_precision}).

\paragraph{Joint variation of threshold and filter size.}
\Cref{fig:operating_grid} varies the probability threshold and the minimum
component size jointly and plots lesion recall against lesion-wise
precision, the two quantities that make up lesion-wise F1. Each faint point
is one setting; the lines connect the settings that are not outperformed
in both precision and recall by another setting of the same loss. On
MSLesSeg, the CATMIL frontier lies above those of DiceCE and CC-DiceCE over
most of the precision range and close to those of the Tversky-based losses,
which is the pattern quantified in \cref{tab:matched_precision}. On 3D-MR-MS, the
default output of CATMIL (filled circle) lies far to the left of its
frontier: its precision of 0.37 is due to the many small components, and
the frontier reaches the precision of the baselines with similar recall.
The minimum component size is given in voxels; a size in mm$^3$
transfers poorly between the datasets: because the
voxels of 3D-MR-MS are 5.7 times smaller, a filter of 15~mm$^3$ removes
85 voxels and with them most of the small lesions. The false-positive
components of CATMIL are therefore small in voxels rather than in
millimeters, and the filter size should be chosen for each image
resolution.

\begin{figure}[htbp]
    \centering
    \includegraphics[width=\textwidth]{figures/operating_grid.pdf}
    \caption{Lesion recall against lesion-wise precision for all
    combinations of probability threshold ($t = 0.1, \dots, 0.9$) and
    minimum component size ($s = 1$ to 170 voxels). Top: small-lesion
    recall ($\leq 150$ voxels); bottom: recall of all lesions. Faint
    points: individual settings; lines: settings not outperformed in both
    precision and recall by another setting of the same loss. Filled
    markers: default output ($t = 0.5$, $s = 1$); open markers: PP5
    ($t = 0.5$, $s = 5$), shown for CATMIL and DiceCE. Metrics are computed
    per case and averaged over cases and the five fold models.}
    \label{fig:operating_grid}
\end{figure}

\begin{figure}[htbp]
    \centering
    \includegraphics[width=\textwidth]{figures/operating_curve.pdf}
    \caption{Lesion recall against false-positive components per case on
    the MSLesSeg test set. Left: predicted components smaller than $s$
    voxels are removed ($s = 0, 2, 3, 5, 8, 10, 15, 20$) at probability
    threshold 0.5. Right: the probability threshold $t$ is varied from 0.1
    to 0.9 without filtering. Open circles: unfiltered predictions at
    $t = 0.5$; stars: PP5 ($s = 5$). A false-positive component is a
    predicted connected component that does not overlap any reference
    lesion. Recall is pooled over the test lesions; values are means over
    the five fold models.}
    \label{fig:operating_curve}
\end{figure}

\paragraph{Lesion-level analyses on the 3D-MR-MS dataset.}
\Cref{fig:recall_vs_size_curve_3dmrms,fig:lesion_fate_3dmrms,fig:operating_curve_3dmrms,fig:max_prob_by_size_3dmrms}
repeat the analyses of
\cref{fig:recall_vs_size_curve,fig:lesion_fate,fig:operating_curve,fig:max_prob_by_size}
on the 3D-MR-MS dataset, and are read in the same way.

The recall curve (\cref{fig:recall_vs_size_curve_3dmrms}) reproduces the
main result. CATMIL detects 51\% of lesions of at most 10~mm$^3$, compared
with 40--41\% for the voxel-level losses and 47\% for CC-DiceCE, and 93\%
of lesions of 11--50~mm$^3$, compared with 86--88\% and 91\%. Small lesions
are far more frequent in this dataset: 182 of the 389 reference lesions are
at most 10~mm$^3$, so the gain in this range accounts for most of the
overall increase in lesion recall, from 0.66--0.67 for the voxel-level
losses and 0.71 for CC-DiceCE to 0.74. Above 100~mm$^3$, CATMIL detects
slightly fewer lesions than the voxel-level losses (95\% against 98--99\%),
a difference of about one of the 34 lesions in this range, and slightly
more than CC-DiceCE (93\%).

The paired comparison (\cref{fig:lesion_fate_3dmrms}) is again one-sided.
CATMIL detects 33.2 lesions per fold model that DiceCE misses, against 2.4
in the opposite direction, 30.4 against 3.8 for each of the Tversky-based
losses, and 17.0 against 3.4 for CC-DiceCE. About 97\% of the lesions
recovered by CATMIL relative to the voxel-level losses, and 89\% of those
recovered relative to CC-DiceCE, are at most 50~mm$^3$.

The operating curve (\cref{fig:operating_curve_3dmrms}) shows that the
component-level cost of CATMIL is larger on this dataset. Without
filtering, CATMIL produces 150 false-positive components per case, compared
with 14--19 for the voxel-level losses and 43 for CC-DiceCE, which explains its low Lesion F1 in
\cref{tab:external_eval_3dmrms}. The size filter again removes most of
them: PP5 reduces the count to 27 per case with a lesion recall (pooled over all reference lesions) of 0.700, and a
filter of 20 voxels reduces it to 9.7 per case with a recall of 0.657. At
this setting, CATMIL has about as many false-positive components as DiceCE
after PP5 (9.5 per case) and a higher recall (0.648 for DiceCE). A larger filter
is therefore needed at the finer in-plane resolution of 3D-MR-MS.

The probability levels (\cref{fig:max_prob_by_size_3dmrms}) show the same
mechanism as on MSLesSeg. For lesions of at most 10~mm$^3$, the fraction
with no response falls from 57--58\% for the voxel-level losses and 51\%
for CC-DiceCE to 39\% for CATMIL.

\begin{figure}[htbp]
    \centering
    \includegraphics[width=0.85\textwidth]{figures/recall_vs_size_curve_3dmrms.pdf}
    \caption{Lesion recall against lesion size on the 3D-MR-MS test set.
    Layout as in \cref{fig:recall_vs_size_curve}.}
    \label{fig:recall_vs_size_curve_3dmrms}
\end{figure}

\begin{figure}[htbp]
    \centering
    \includegraphics[width=\textwidth]{figures/lesion_fate_3dmrms.pdf}
    \caption{Lesions detected by only one of two losses on the 3D-MR-MS
    test set, by lesion size. Layout as in \cref{fig:lesion_fate}.}
    \label{fig:lesion_fate_3dmrms}
\end{figure}

\begin{figure}[htbp]
    \centering
    \includegraphics[width=\textwidth]{figures/operating_curve_3dmrms.pdf}
    \caption{Lesion recall against false-positive components per case on
    the 3D-MR-MS test set. Layout as in \cref{fig:operating_curve}.}
    \label{fig:operating_curve_3dmrms}
\end{figure}

\begin{figure}[htbp]
    \centering
    \includegraphics[width=0.75\textwidth]{figures/max_prob_by_size_3dmrms.pdf}
    \caption{Maximum foreground probability inside each reference lesion on
    the 3D-MR-MS test set, by lesion size. Layout as in
    \cref{fig:max_prob_by_size}.}
    \label{fig:max_prob_by_size_3dmrms}
\end{figure}

\section{Robustness to the Detection Criterion}
\label{appendix:strict-detection}

All lesion-level results in the main text count a reference lesion as
detected if any predicted voxel overlaps it, and the case-based analysis
shows that CATMIL often covers only part of the lesions it recovers. \Cref{tab:strict_detection} therefore recomputes
small-lesion recall with stricter criteria, which require a minimum
fraction of the lesion voxels to be covered by the default segmentation
($p > 0.5$).

\begin{table}[htbp]
\centering
\caption{Small-lesion recall ($\leq 150$ voxels) under detection criteria
of increasing strictness. A lesion counts as detected if the prediction
covers any of its voxels, or at least the stated fraction of them. Recall
is pooled over the reference lesions of the test set and averaged over the
five fold models, so the any-overlap values differ slightly from the
case-averaged Small Lesion Recall in the main tables. $\Delta$ and bold
as in \cref{tab:clinical_size}.}
\label{tab:strict_detection}
\resizebox{\textwidth}{!}{%
\begin{tabular}{llcccccll}
\toprule
\textbf{Dataset} & \textbf{Criterion} & \textbf{CATMIL} & \textbf{DiceCE} & \textbf{Tversky} & \textbf{FocalTversky} & \textbf{CC-DiceCE} & $\Delta$ vs.\ DiceCE & $\Delta$ vs.\ CC-DiceCE \\
\midrule
MSLesSeg & any overlap & 0.883 & 0.779 & 0.829 & 0.833 & 0.841 & \textbf{+0.104 [+0.046, +0.176]} & \textbf{+0.042 [+0.024, +0.062]} \\
 & $\geq$ 10\% & 0.865 & 0.767 & 0.818 & 0.818 & 0.826 & \textbf{+0.098 [+0.039, +0.170]} & \textbf{+0.039 [+0.020, +0.060]} \\
 & $\geq$ 25\% & 0.836 & 0.743 & 0.802 & 0.799 & 0.799 & \textbf{+0.093 [+0.032, +0.162]} & \textbf{+0.037 [+0.012, +0.062]} \\
 & $\geq$ 50\% & 0.760 & 0.679 & 0.750 & 0.748 & 0.725 & \textbf{+0.081 [+0.029, +0.138]} & \textbf{+0.035 [+0.007, +0.066]} \\
\midrule
3D-MR-MS & any overlap & 0.654 & 0.548 & 0.564 & 0.562 & 0.614 & \textbf{+0.107 [+0.046, +0.129]} & \textbf{+0.041 [+0.011, +0.080]} \\
 & $\geq$ 10\% & 0.630 & 0.538 & 0.555 & 0.556 & 0.602 & \textbf{+0.092 [+0.029, +0.114]} & +0.028 [$-$0.004, +0.062] \\
 & $\geq$ 25\% & 0.595 & 0.509 & 0.534 & 0.538 & 0.584 & \textbf{+0.086 [+0.009, +0.108]} & +0.011 [$-$0.034, +0.046] \\
 & $\geq$ 50\% & 0.550 & 0.449 & 0.493 & 0.500 & 0.535 & \textbf{+0.101 [+0.028, +0.129]} & \textbf{+0.015 [+0.002, +0.046]} \\
\bottomrule
\end{tabular}%
}
\end{table}

The gain over DiceCE does not depend on the lenient criterion. On both
datasets it stays between 8 and 11 percentage points at every coverage
level, and its CI excludes zero in every case. The margin over the
Tversky-based losses narrows as the criterion becomes stricter: at 50\%
coverage on MSLesSeg, CATMIL detects 76.0\% of small lesions against
74.8--75.0\%. The gain over CC-DiceCE stays between 3.5 and 4.2 points at
every coverage level on MSLesSeg, with CIs that exclude zero, but on
3D-MR-MS it shrinks from 4.1 points under any overlap to 1.1--2.8 points
under the coverage criteria, with CIs that include zero at 10\% and 25\%.
This agrees with the case-based analysis: the lesion-level term ensures that a lesion is detected, whereas its
full extent still depends on the voxel-wise terms. Lesion-wise precision
cannot be recomputed with a coverage criterion from the stored lesion
tables, so this analysis is limited to recall.

\section{Detection at Matched Lesion Precision}
\label{sec:matched_precision}

PP5 compares the losses at one fixed setting, at which CATMIL still has a
lower lesion-wise precision than the baselines. To test whether CATMIL
detects more lesions rather than only operating at a more liberal point,
each loss was evaluated over a grid of probability thresholds
($t = 0.1, \dots, 0.9$) and minimum component sizes ($s = 1$ to 170
voxels) (\cref{fig:operating_grid} in \cref{appendix:lesion-level-analyses}).
For each baseline, the CATMIL setting with the highest small-lesion recall
whose precision is at least that of the baseline at its default output was
selected on MSLesSeg and applied unchanged to 3D-MR-MS
(\cref{tab:matched_precision}).

\begin{table}[t]
\centering
\caption{Small-lesion recall and lesion-wise F1 of CATMIL at matched
lesion-wise precision (setting selected on MSLesSeg and applied unchanged
to 3D-MR-MS; the MSLesSeg values are therefore optimistic). Values are
CATMIL minus baseline with 95\% patient-level bootstrap CIs in brackets;
bold: CI excludes zero.}
\label{tab:matched_precision}
\resizebox{\textwidth}{!}{%
\begin{tabular}{l r@{\hspace{0.4em}}l r@{\hspace{0.4em}}l r@{\hspace{0.4em}}l r@{\hspace{0.4em}}l}
\toprule
& \multicolumn{2}{c}{\textbf{vs.\ DiceCE}} & \multicolumn{2}{c}{\textbf{vs.\ Tversky}} & \multicolumn{2}{c}{\textbf{vs.\ FocalTversky}} & \multicolumn{2}{c}{\textbf{vs.\ CC-DiceCE}} \\
\cmidrule(lr){2-3}\cmidrule(lr){4-5}\cmidrule(lr){6-7}\cmidrule(lr){8-9}
CATMIL $(t, s)$ & \multicolumn{2}{c}{(0.1, 15)} & \multicolumn{2}{c}{(0.8, 4)} & \multicolumn{2}{c}{(0.2, 10)} & \multicolumn{2}{c}{(0.1, 5)} \\
\midrule
\multicolumn{9}{l}{\textit{MSLesSeg (selection)}} \\
\quad $\Delta$ Small Recall & \textbf{+0.048} & [+0.013, +0.111] & +0.016 & [$-$0.013, +0.060] & +0.016 & [$-$0.012, +0.057] & \textbf{+0.047} & [+0.020, +0.084] \\
\quad $\Delta$ Lesion F1    & \textbf{+0.021} & [+0.005, +0.041] & +0.007 & [$-$0.017, +0.028] & +0.007 & [$-$0.008, +0.023] & +0.015 & [$-$0.004, +0.034] \\
\midrule
\multicolumn{9}{l}{\textit{3D-MR-MS (transfer)}} \\
\quad $\Delta$ Small Recall & +0.040 & [$-$0.002, +0.087] & $-$0.010 & [$-$0.049, +0.024] & +0.006 & [$-$0.036, +0.043] & \textbf{+0.039} & [+0.003, +0.080] \\
\quad $\Delta$ Precision    & $-$0.050 & [$-$0.099, +0.001] & +0.030 & [$-$0.040, +0.094] & +0.006 & [$-$0.050, +0.067] & \textbf{$-$0.075} & [$-$0.129, $-$0.026] \\
\quad $\Delta$ Lesion F1    & $-$0.008 & [$-$0.036, +0.023] & +0.008 & [$-$0.020, +0.037] & +0.007 & [$-$0.022, +0.040] & $-$0.031 & [$-$0.072, +0.006] \\
\bottomrule
\end{tabular}%
}
\end{table}

\textbf{At the precision of DiceCE, CATMIL detects more small lesions and
has a higher lesion-wise F1.} With $t = 0.1$ and $s = 15$, CATMIL reaches
the precision of DiceCE on MSLesSeg (0.853 against 0.852) with a
small-lesion recall that is 0.048 higher and a lesion-wise F1 that is
0.021 higher, both supported in all six patients. Transferred unchanged to
3D-MR-MS, the recall gain is of similar size (+0.040), but its CI narrowly
includes zero, and the precision is 0.050 lower.

\textbf{At the precision of CC-DiceCE, CATMIL still detects more small
lesions.} A filter of 5 voxels at $t = 0.1$ brings CATMIL to the precision
of CC-DiceCE on MSLesSeg (0.767 against 0.765) with a small-lesion recall
that is 0.047 higher in all six patients; the lesion-wise F1 difference
(+0.015) is not supported. On 3D-MR-MS, the recall gain is of the same size
(+0.039 [+0.003, +0.080]), but the precision is 0.075 lower, so on that
dataset the gain comes partly from a more sensitive operating point.

\textbf{At matched precision, CATMIL does not detect more lesions than the
Tversky-based losses.} The gain shrinks to +0.016 on MSLesSeg and to
between $-0.010$ and +0.006 on 3D-MR-MS, with CIs that include zero. When
every loss is tuned to its own maximum lesion-wise F1, the best F1 on
MSLesSeg is 0.876 for CATMIL and 0.859--0.872 for the baselines, and none
of the differences is supported.

\section{Effect of the Architecture and of the Objective}
\label{sec:arch_vs_objective}

To test whether a different network could recover the missed small lesions
without changing the objective, four further architectures trained with
DiceCE were evaluated on the MSLesSeg test set under the same protocol: the
residual convolutional network SegResNet~\cite{crimi_3d_2019}, the
transformer-based SwinUNETR~\cite{hatamizadeh_swin_2022} and
UNETR~\cite{hatamizadeh_unetr_2021}, and the state-space model U-Mamba in
its UMambaBot variant~\cite{ma_u-mamba_2024}. They were trained with the
same five cross-validation splits and 150-epoch budget as the nnU-Net
models, with each architecture's default hyperparameters
(\cref{tab:arch_vs_objective}).

\begin{table}[t]
\centering
\caption{Small-lesion detection by architecture and objective on the
MSLesSeg test set, without post-processing. All architectures are trained
with the DiceCE objective, except nnU-Net with CATMIL. Values are means
over the test scans and the five fold models; formatting as in
\cref{tab:loss_comparison_small_lesion}.}
\label{tab:arch_vs_objective}
\resizebox{\textwidth}{!}{%
\begin{tabular}{lcccccc}
\toprule
\textbf{Metric}
& \textbf{nnU-Net}
& \textbf{nnU-Net}
& \textbf{SegResNet}
& \textbf{UMambaBot}
& \textbf{SwinUNETR}
& \textbf{UNETR} \\
& \textbf{(CATMIL)}
& \textbf{(DiceCE)}
& \textbf{(DiceCE)}
& \textbf{(DiceCE)}
& \textbf{(DiceCE)}
& \textbf{(DiceCE)} \\
\midrule
Small Lesion Recall $\uparrow$
& \textbf{0.873} & \underline{0.796} & 0.779 & 0.778 & 0.771 & 0.715 \\
FN Count $\downarrow$
& \textbf{2.57} & 4.95 & 4.83 & \underline{4.65} & 4.92 & 6.43 \\
Lesion F1 $\uparrow$
& 0.757 & \textbf{0.843} & 0.836 & \underline{0.842} & 0.804 & 0.765 \\
\midrule
DSC $\uparrow$
& \textbf{0.783} & 0.780 & 0.761 & \underline{0.782} & 0.762 & 0.745 \\
HD95 (mm) $\downarrow$
& \textbf{7.98} & 9.04 & 10.58 & \underline{8.68} & 10.57 & 10.96 \\
\bottomrule
\end{tabular}%
}
\end{table}

\textbf{Changing the architecture did not recover the missed small
lesions.} Trained with DiceCE, the four architectures reached a
small-lesion recall similar to or below that of nnU-Net (0.715--0.779,
against 0.796); the paired difference was below zero for UNETR ($-0.081$
[$-0.170$, $-0.021$]) and included zero for the others.

\textbf{Changing the objective recovered them.} With CATMIL, nnU-Net
exceeded every architecture trained with DiceCE in all six patients, by
+0.094 [+0.022, +0.206] (SegResNet) to +0.158 [+0.067, +0.302] (UNETR), and
missed 2.08 to 3.87 fewer lesions per scan, all CIs excluding zero; the
gain remained after PP5 (+0.091 to +0.168). UMambaBot matched the DSC of
nnU-Net with DiceCE (0.782 against 0.780) without detecting more small
lesions (0.778), so an overlap score does not reveal this difference.
Whether CATMIL raises small-lesion recall as much when combined with these
architectures was not tested.

\textbf{The result is reproduced on the cross-validation patients.} On
the cross-validation predictions of the 47 MSLesSeg training patients
(\cref{sec:catmil_oof}), CATMIL exceeded each architecture trained with
DiceCE by 0.100 to 0.146 in small-lesion recall
(\cref{tab:oof_arch}).

\begin{table}[t]
\centering
\caption{Paired differences between nnU-Net with CATMIL and four further
architectures trained with DiceCE on the cross-validation predictions of
the 81 MSLesSeg training scans (47 patients). Values and formatting as in
\cref{tab:oof_paired}.}
\label{tab:oof_arch}
\resizebox{\textwidth}{!}{%
\begin{tabular}{lccccc}
\toprule
\textbf{CATMIL minus} & \textbf{Small Lesion Recall} & \textbf{FN Count} & \textbf{Lesion F1} & \textbf{Lesion F1 (PP5)} & \textbf{DSC} \\
\midrule
SegResNet (DiceCE) & \textbf{+0.107 [+0.076, +0.136]} & \textbf{$-$3.12 [$-$4.46, $-$2.10]} & \textbf{$-$0.098 [$-$0.131, $-$0.064]} & $-$0.018 [$-$0.037, +0.002] & +0.000 [$-$0.008, +0.009] \\
UMambaBot (DiceCE) & \textbf{+0.112 [+0.075, +0.144]} & \textbf{$-$3.53 [$-$4.96, $-$2.41]} & \textbf{$-$0.106 [$-$0.140, $-$0.072]} & \textbf{$-$0.030 [$-$0.051, $-$0.010]} & +0.001 [$-$0.008, +0.009] \\
SwinUNETR (DiceCE) & \textbf{+0.100 [+0.065, +0.132]} & \textbf{$-$3.28 [$-$4.80, $-$2.12]} & \textbf{$-$0.058 [$-$0.099, $-$0.015]} & +0.007 [$-$0.018, +0.032] & \textbf{+0.028 [+0.014, +0.044]} \\
UNETR (DiceCE)     & \textbf{+0.146 [+0.105, +0.187]} & \textbf{$-$4.59 [$-$6.49, $-$3.06]} & $-$0.009 [$-$0.055, +0.040] & \textbf{+0.057 [+0.026, +0.091]} & \textbf{+0.051 [+0.032, +0.074]} \\
\bottomrule
\end{tabular}%
}
\end{table}

\section{Gradient and Boundedness Properties of the CATMIL Objective}
\label{appendix:catmil-properties}

This section derives the properties of the CAT and MIL terms stated in
\cref{sec:method}, using the notation of that section. The components $C_k$
and the weights $w_i$ are computed from the ground-truth mask and are
constants with respect to the predictions $p_i$.

\paragraph{Complete definition of the CAT term.}
In general form, the foreground weight is
$w_i = (|C_k| + \varepsilon_{c})^{-\gamma}$ for $i \in C_k$, where the
exponent $\gamma \geq 0$ controls the degree of size compensation:
$\gamma=0$ keeps the weighting proportional to the number of voxels, and
$\gamma=1$ is used in this work (\cref{eq:sl_cat_weights}). The offset
$\varepsilon_{c}=5$ voxels limits the per-voxel weight of very small lesions,
and $\rho_k \in [1/6, 1)$ for any lesion of at least one voxel. Because the foreground weights are much
less than one, the background weight is set to $w_{\mathrm{bg}}=0.01$ so
that the false-positive term remains on a comparable scale. The weighted
sums are
\begin{equation}
\mathrm{TP}_{w} = \sum_{i \in \Omega} w_i\, p_i\, g_i,
\qquad
\mathrm{FP}_{w} = \sum_{i \in \Omega} w_i\, p_i\, (1-g_i),
\qquad
\mathrm{FN}_{w} = \sum_{i \in \Omega} w_i\, (1-p_i)\, g_i ,
\end{equation}
and $\delta=10^{-5}$ in \cref{eq:sl_cat_loss} is the usual smoothing
constant. Both the CAT and the MIL term are set to zero when the training
batch contains no lesions. The base loss retains nnU-Net's deep
supervision, whereas the CAT and MIL terms are computed only on the
full-resolution output.

\paragraph{Lesion-level form of the CAT term.}
Let $r_k = \frac{1}{|C_k|}\sum_{i \in C_k} p_i \in [0,1]$ be the mean
predicted probability inside lesion $C_k$, that is, its soft recall. With
$\gamma=1$,
\begin{equation}
\mathrm{TP}_{w} = \sum_{k=1}^{K} \rho_k\, r_k,
\qquad
\mathrm{FN}_{w} = \sum_{k=1}^{K} \rho_k\, (1-r_k) .
\label{eq:sl_cat_lesion_form}
\end{equation}
Lesions therefore contribute to the loss through their recalls $r_k$, with
coefficients $\rho_k$ that are nearly independent of lesion size. A
completely missed lesion ($r_m=0$) transfers a total weight of
$\rho_m < 1$ from $\mathrm{TP}_{w}$ to $\mathrm{FN}_{w}$, whereas its
effect on the DSC is proportional to its volume $|C_m|$: if a prediction is
perfect except that lesion $C_m$ is missed,
\begin{equation}
\mathrm{DSC}
=
\frac{2\,(|G| - |C_m|)}{2|G| - |C_m|}
=
1 - \frac{|C_m|}{2|G| - |C_m|} .
\label{eq:sl_missed_lesion}
\end{equation}
For $\gamma=0$ and $w_{\mathrm{bg}}=1$, the CAT term reduces to the
standard Tversky loss; if additionally $\alpha=\beta=0.5$, it is the soft
Dice loss.

\paragraph{Gradient of the CAT term.}
Let
\begin{equation}
A = \mathrm{TP}_{w}+\delta,
\qquad
B = A + \alpha\, \mathrm{FP}_{w} + \beta\, \mathrm{FN}_{w},
\end{equation}
so that $\mathcal{L}_{\mathrm{CAT}} = 1 - A/B$. For a foreground voxel
$j \in C_k$, $\partial \mathrm{TP}_{w}/\partial p_j = w_j$,
$\partial \mathrm{FP}_{w}/\partial p_j = 0$, and
$\partial \mathrm{FN}_{w}/\partial p_j = -w_j$, hence
$\partial A/\partial p_j = w_j$ and
$\partial B/\partial p_j = (1-\beta)\, w_j$. The quotient rule gives
\begin{equation}
\frac{\partial \mathcal{L}_{\mathrm{CAT}}}{\partial p_j}
=
-\frac{w_j B - (1-\beta)\, w_j A}{B^2}
=
-c\, w_j ,
\qquad
c
=
\frac{\beta A + \alpha\, \mathrm{FP}_{w} + \beta\, \mathrm{FN}_{w}}{B^2}
> 0 .
\end{equation}
The factor $c$ is the same for all foreground voxels, so the gradient at a
lesion voxel is scaled only by its weight
$w_j = (|C_k|+\varepsilon_{c})^{-\gamma}$. Summing over the component, the
total gradient magnitude received by lesion $C_k$ is
\begin{equation}
\sum_{j \in C_k}
\left| \frac{\partial \mathcal{L}_{\mathrm{CAT}}}{\partial p_j} \right|
=
c\, |C_k|\, (|C_k|+\varepsilon_{c})^{-\gamma},
\end{equation}
which is proportional to $|C_k|$ for $\gamma=0$ and equal to
$c\,\rho_k < c$ for $\gamma=1$, with $\rho_k$ the lesion weight of
\cref{eq:sl_cat_lesion_weight}.

For a background voxel $j$ with $g_j=0$, only $\mathrm{FP}_{w}$ depends on
$p_j$, so $\partial A/\partial p_j = 0$ and
$\partial B/\partial p_j = \alpha\, w_{\mathrm{bg}}$, and
\begin{equation}
\frac{\partial \mathcal{L}_{\mathrm{CAT}}}{\partial p_j}
=
\frac{\alpha\, w_{\mathrm{bg}}\, A}{B^2}
\geq 0 .
\end{equation}
Increasing the predicted probability at a background voxel therefore
increases the loss.

\paragraph{Gradient of the MIL term.}
With the clipping of \cref{eq:sl_mil_loss}, the contribution of lesion
$C_k$ is $-\log\bigl(\max(s_k,\varepsilon)\bigr)$ and remains finite. Let
$m_k = \arg\max_{i \in C_k} p_i$ be the maximal voxel of lesion $C_k$.
Away from ties,
\begin{equation}
\frac{\partial \mathcal{L}_{\mathrm{MIL}}}{\partial p_j}
=
\begin{cases}
-\dfrac{1}{K s_k},
& \text{if } j=m_k \text{ for some } k \text{ and } s_k>\varepsilon, \\[2mm]
0, & \text{otherwise.}
\end{cases}
\end{equation}
The gradient is concentrated on one voxel per lesion, its magnitude does
not depend on $|C_k|$, and it grows as the strongest response $s_k$
decreases. When several voxels share the maximum, the hard maximum is not
differentiable but remains subdifferentiable, and automatic differentiation
assigns the gradient to a maximal voxel, as in max-pooling.

\paragraph{Boundedness and continuity.}
Since $\mathrm{TP}_{w}$, $\mathrm{FP}_{w}$, $\mathrm{FN}_{w}$, $\alpha$, and
$\beta$ are non-negative and $\delta>0$, $0 < A \leq B$, and therefore
$0 \leq \mathcal{L}_{\mathrm{CAT}} < 1$. Since
$\varepsilon \leq \max(s_k,\varepsilon) \leq 1$,
$0 \leq \mathcal{L}_{\mathrm{MIL}} \leq -\log \varepsilon$. For
non-negative coefficients the auxiliary part of the objective thus
satisfies
\begin{equation}
0
\leq
\lambda_{\mathrm{CAT}}\, \mathcal{L}_{\mathrm{CAT}}
+
\lambda_{\mathrm{MIL}}\, \mathcal{L}_{\mathrm{MIL}}
\leq
\lambda_{\mathrm{CAT}} - \lambda_{\mathrm{MIL}} \log \varepsilon .
\end{equation}
Both terms are continuous in the predictions. The CAT term is a ratio of
linear functions with a strictly positive denominator and is differentiable
everywhere. The MIL term is differentiable except at ties of the maximum
and at the clipping point $s_k=\varepsilon$, which is the same piecewise
behavior as that of max-pooling and does not prevent gradient-based
training.
